% TEMPLATE-MILC-v3.tex - plantilla maestra UNICA de las guias MILC v3.
%
% La usan las cuatro familias de guias, cada una con su driver:
%   · Grados 6-11          -> scripts/build-guias-g11.py
%   · Bebras               -> scripts/build-guias-bebras.py
%   · Semillero            -> scripts/build-guias-semillero.py
%   · Territorio interior  -> scripts/build-guias-territorio-interior.py
%
% Antes habia tres copias casi identicas (~400 lineas iguales, 6 distintas):
% cada mejora habia que aplicarla tres veces, y en la practica no se hacia
% (el motor de recursos vivia solo en dos de ellas). Lo que varia entre
% familias esta parametrizado en 6 placeholders que cada driver llena
% (se nombran aqui SIN su sintaxis real: el builder reemplaza por texto plano,
% tambien dentro de los comentarios, y un valor multilinea rompe el preambulo):
%
%   HEADER_IZQ        encabezado izquierdo de pagina
%   HEADER_DER        encabezado derecho de pagina
%   PORTADA_ETIQUETA  rotulo sobre el numero grande (GRADO/MOMENTO/MODULO)
%   PORTADA_SERIE     linea de serie de la portada
%   PORTADA_ID        identificador de la guia en la portada
%   PORTADA_CIRCULO   el circulito de grado (°), o vacio si no aplica
%   PDF_TITULO        titulo en texto plano (sin \textbf ni comillas TeX) para
%                     los metadatos del PDF (/Title); lo produce lib_xelatex.texto_pdf
%
% Composicion (auditoria 2026-09-03): las bandas de estacion piden espacio
% (\Needspace*) y prohiben el salto justo despues (\nopagebreak), asi no quedan
% huerfanas al pie; las cajas largas (softbox, drawbox, infoband, puentebox) son
% `breakable`, asi una caja de media pagina ya no arrastra todo a la siguiente
% dejando la anterior al 40 %. Todo texto sobre mostaza o verde va en milcNegro
% (blanco sobre #E5B400 da 1,93:1 y sobre #58A923 2,95:1; ninguno pasa WCAG AA).
% El resto de claves las documenta content/guias/_SCHEMA.md. El script valida
% que no queden placeholders sin reemplazar antes de escribir el archivo final.

% Etiquetado PDF (latex-lab): /Lang, estructura y texto alternativo de las
% figuras, para lectores de pantalla. Debe ir ANTES de \documentclass.
\DocumentMetadata{lang=es-CO,tagging=on}
\documentclass[11pt,letterpaper]{article}
% headheight=24pt: la cabecera derecha lleva dos lineas (identidad de la guia
% y estacion actual). headsep se acorta en la misma medida para que el bloque
% de texto quede exactamente donde estaba con headheight=12pt.
\usepackage[margin=2.0cm,headheight=24pt,headsep=13pt]{geometry}
\usepackage[table]{xcolor}
\usepackage{fontspec}
\usepackage[spanish]{babel}
\usepackage{tikz}
% Librerías TikZ para los diagramas de las guías (bloque `recursos:`):
%  · babel  → babel en español vuelve activos caracteres como «>», y TikZ los usa
%             en sus opciones (>=stealth, ->). Sin esta librería, cualquier
%             diagrama con flechas revienta con «\language@active@arg> has an
%             extra }». Debe ir ANTES de que se dibuje nada.
%  · positioning        → sintaxis «below left=2pt and 3pt».
%  · decorations.pathreplacing → llaves ({brace}) para señalar rangos.
\usetikzlibrary{babel,positioning,decorations.pathreplacing}
\usepackage{graphicx}
% Circuitos: el trabajo de electrónica se hace hoy con micro:bit y sus sensores
% integrados (MakeCode, sin cableado), así que no se necesitan esquemáticos.
% Si algún día entran montajes con protoboard, basta descomentar esta línea:
% los diagramas .tex/.tikz declarados en `recursos:` ya se incrustan solos.
% \usepackage{circuitikz}
\usepackage[most]{tcolorbox}
\usepackage{tabularx}
\usepackage{array}
\usepackage{fancyhdr}
\usepackage{titlesec}
\usepackage[document]{ragged2e}  % bandera a la derecha en todo el cuerpo
\usepackage{enumitem}
\usepackage{microtype}
\usepackage{etoolbox}
\usepackage{needspace}
% hyperref al final del preambulo: metadatos del PDF (titulo, autor, idioma,
% familia), marcadores por estacion (\pdfbookmark) y enlaces sin recuadro.
\usepackage[hidelinks,
  pdftitle={Mentoría: acompañar a alguien que empieza donde tú empezaste},
  pdfauthor={PhD. Álvaro Cárdenas Orozco},
  pdfsubject={Territorio interior · Educación socioemocional · Grado 11° · Momento 9 · MILC v3},
  pdflang={es-CO},
  bookmarksopen=true,bookmarksnumbered=false]{hyperref}

% Helvetica Neue no trae la flecha U+2192, asi que cada «→» escrito literalmente
% en el YAML se compilaba a NADA, en silencio (462 flechas en 61 guias: rutas del
% tipo «paleta → tipografia → lockup» salian rotas). Mapeamos el caracter a su
% version matematica, que si tiene glifo. Asi el autor puede seguir escribiendo
% «→» comodamente en el YAML.
\usepackage{newunicodechar}
\newunicodechar{→}{\ensuremath{\rightarrow}}
% Mismo problema con los simbolos de marca y vinetas: el «✓» de «una palomita ✓»
% salia como cuadrito vacio en el PDF de todas las guias que lo usaban.
\usepackage{pifont}
\newunicodechar{✓}{\ding{51}}
\newunicodechar{✗}{\ding{55}}
\newunicodechar{✔}{\ding{52}}
\newunicodechar{•}{\textbullet}
\newunicodechar{≥}{\ensuremath{\geq}}
\newunicodechar{≤}{\ensuremath{\leq}}

\setmainfont{Helvetica Neue}
\setsansfont{Helvetica Neue}
\setlength{\parindent}{0pt}
\setlength{\parskip}{6pt}
% Sin particion de palabras: en 6.º-7.º una palabra cortada por guion al final
% de linea («Comuni-cacion») frena la lectura mucho mas que una linea corta.
\hyphenpenalty=10000
\exhyphenpenalty=10000
\pretolerance=2000
\tolerance=3000
\setlength{\emergencystretch}{3em}
\linespread{1.10}
\renewcommand{\arraystretch}{1.25}
\setlist{leftmargin=5mm,itemsep=1mm,topsep=1mm}
\newcolumntype{Y}{>{\RaggedRight\arraybackslash}X}

\definecolor{milcMagenta}{HTML}{E6007E}
\definecolor{milcVino}{HTML}{5A0038}
\definecolor{milcTurquesa}{HTML}{0093A5}
\definecolor{milcVerde}{HTML}{58A923}
\definecolor{milcMostaza}{HTML}{E5B400}
\definecolor{milcOcre}{HTML}{B86600}
\definecolor{milcNegro}{HTML}{241816}
\definecolor{milcGris}{HTML}{F7F7F4}
\definecolor{milcLinea}{HTML}{E8E2DD}
\definecolor{milcGradePrimary}{HTML}{0D9488}
\definecolor{milcGradeDark}{HTML}{115E59}
\definecolor{milcGradeSoft}{HTML}{CCFBF1}
\definecolor{milcGradeAccent}{HTML}{CFE7FF}
\definecolor{milcGradeText}{HTML}{FFFFFF}
\color{milcNegro}

\pagestyle{fancy}
\fancyhf{}
% Cabecera de dos lineas: identidad de la guia arriba y, a la derecha y en
% gris, la estacion en curso (\markright desde \milcestacion). Antes de la
% primera estacion la segunda linea va vacia; el \strut mantiene la altura
% para que la linea de arriba no baile entre paginas.
\lhead{\small Territorio interior · Educación socioemocional\\[-1pt]{\footnotesize\strut}}
\rhead{\small Grado 11° · Momento 9 · MILC v3\\[-1pt]{\footnotesize\strut\color{milcNegro!65}\rightmark}}
\cfoot{\small\thepage}
\renewcommand{\headrulewidth}{0pt}

% Sobre mostaza (#E5B400) y verde (#58A923) el texto va en milcNegro; sobre el
% resto de la paleta, en blanco. Se usa dentro de fonttitle/fontupper, donde un
% \color posterior manda sobre coltitle.
\newcommand{\milctintasobre}[1]{%
  \ifboolexpr{test{\ifstrequal{#1}{milcMostaza}} or test{\ifstrequal{#1}{milcVerde}}}%
    {\color{milcNegro}}{\color{white}}}

\newtcolorbox{titlebox}[1]{
  enhanced,colback=#1,colframe=#1,arc=7mm,boxrule=0pt,
  left=7mm,right=7mm,top=3mm,bottom=3mm,
  before skip=8pt,after skip=8pt,
  fontupper=\sffamily\bfseries\Large\color{white}
}

\newtcolorbox{softbox}[3]{
  enhanced,breakable,pad at break=3mm,
  colback=#2,colframe=#1,borderline west={4pt}{0pt}{#1},
  arc=2mm,boxrule=.4pt,left=4mm,right=4mm,top=3mm,bottom=3mm,
  before skip=6pt,after skip=8pt,title={#3},coltitle=white,
  colbacktitle=#1,fonttitle=\sffamily\bfseries\milctintasobre{#1},fontupper=\RaggedRight,
  attach boxed title to top left={xshift=3mm,yshift=-2mm},
  boxed title style={arc=2mm, boxrule=0pt}
}

\newtcolorbox{drawbox}[2]{
  enhanced,breakable,pad at break=4mm,
  colback=white,colframe=#1,arc=4mm,boxrule=1pt,
  left=5mm,right=5mm,top=4mm,bottom=4mm,before skip=8pt,after skip=8pt,
  title={#2},coltitle=white,colbacktitle=#1,fonttitle=\sffamily\bfseries\milctintasobre{#1},
  fontupper=\RaggedRight,
  attach boxed title to top left={xshift=4mm,yshift=-2mm},
  boxed title style={arc=3mm, boxrule=0pt}
}

\newtcolorbox{infoband}[1]{
  enhanced,breakable,pad at break=2mm,colback=#1!8,colframe=#1!50,arc=2mm,boxrule=.5pt,
  left=4mm,right=4mm,top=2mm,bottom=2mm,before skip=4pt,after skip=4pt,
  fontupper=\small\RaggedRight
}

% Puente narrativo entre secciones (contrato editorial Conéctate).
% Da continuidad: "Acabas de X. Ahora vas a Y porque Z."
% Visualmente: banda izquierda en color del grado, texto en itálica pequeña.
\newtcolorbox{puentebox}{
  enhanced,breakable,pad at break=2mm,colback=milcGradeSoft,colframe=milcGradeDark,
  borderline west={3pt}{0pt}{milcGradeDark},
  arc=0mm,boxrule=0pt,left=5mm,right=4mm,top=2.5mm,bottom=2.5mm,
  before skip=4pt,after skip=8pt,
  fontupper=\itshape\small\color{milcNegro}\RaggedRight
}
% La flecha va en modo matematico a proposito: Helvetica Neue NO tiene el glifo
% U+2192, asi que \textrightarrow se compilaba a NADA («Missing character: There
% is no → in font Helvetica Neue») y el puente salia sin flecha. En modo
% matematico la flecha viene de la fuente matematica y si se dibuja.
\newcommand{\puente}[1]{%
  \begin{puentebox}%
    \textcolor{milcGradeDark}{$\rightarrow$}\hspace{2mm}#1%
  \end{puentebox}%
}

\newcommand{\linead}[1]{\noindent\color{milcLinea}\rule{#1}{0.5pt}\color{milcNegro}}
\newcommand{\respuesta}[1]{\par\vspace{2mm}\foreach \n in {1,...,#1}{\noindent\color{milcLinea}\rule{\linewidth}{0.5pt}\par\vspace{5mm}}\color{milcNegro}}
\newcommand{\checkbox}{\textcolor{milcGradeDark}{\fbox{\phantom{x}}}\hspace{5pt}}

% ─── Listas aireadas para las guías socioemocionales ────────────────────
% Componer las verificaciones como «\checkbox texto\\» deja el texto que salta
% de línea debajo del cuadrito y apelmaza el bloque. Estos entornos alinean la
% continuación y airean los ítems. Los usa Territorio Interior; cualquier
% familia puede hacerlo.
\newlist{checklista}{itemize}{1}
\setlist[checklista]{
  label=\textcolor{milcGradeDark}{\fbox{\phantom{x}}},
  leftmargin=9mm, labelsep=3.2mm, itemsep=2.4mm,
  topsep=2.5mm, parsep=0pt, partopsep=0pt, before=\RaggedRight
}
\newlist{pasoslista}{enumerate}{1}
\setlist[pasoslista]{
  label=\textcolor{milcGradeDark}{\sffamily\bfseries\arabic*.},
  leftmargin=8mm, labelsep=3mm, itemsep=2.8mm,
  topsep=1mm, parsep=0pt, partopsep=0pt, before=\RaggedRight
}

\newcommand{\phasecard}[4]{
\begin{tcolorbox}[enhanced,colback=#3,colframe=#2,arc=3mm,boxrule=.5pt,height=3.05cm,valign=center,halign=center,left=1.5mm,right=1.5mm,top=2mm,bottom=2mm]
{\sffamily\bfseries\fontsize{22}{24}\selectfont\textcolor{#2}{#1}}\par
{\sffamily\bfseries\small #4}
\end{tcolorbox}
}

% ── Piezas del rediseno 2026 (legibilidad 6.º-11.º) ───────────────────
% Banda de estacion: reemplaza el titlebox macizo. Titulo a la izquierda,
% dato practico (tiempo/modalidad) a la derecha, en una sola linea.
% Nunca huerfana: pide 9 lineas libres (o salta de pagina rellenando con
% \vfill) y prohibe el salto justo despues de la banda. Nueve y no seis:
% la caja partible que sigue necesita ~80pt (titulo adosado, relleno y dos
% lineas) para arrancar en la misma pagina; con menos, tcolorbox la manda
% entera a la siguiente y la banda se queda sola. El penalty 10000 va
% en `after`, ANTES del espacio posterior: un `after skip` (aunque sea 0pt)
% es un \vskip tras la caja, y ese glue es un punto de corte legal que un
% \nopagebreak escrito despues de \end{tcolorbox} ya no cierra. Cada banda
% deja marcador PDF y actualiza la cabecera derecha.
\newcounter{milcestacion}
\newcommand{\milcestacion}[2]{%
\Needspace*{9\baselineskip}%
\stepcounter{milcestacion}%
\pdfbookmark[1]{#2}{milcestacion\themilcestacion}%
\markright{#2}%
\begin{tcolorbox}[enhanced,colback=#1,colframe=#1,arc=2mm,boxrule=0pt,
  left=5mm,right=5mm,top=2.6mm,bottom=2.6mm,before skip=11pt,
  after={\par\nopagebreak\vskip7pt}]
{\sffamily\bfseries\fontsize{15}{18}\selectfont\milctintasobre{#1}#2}
\end{tcolorbox}}

% Tarjeta de paso numerada: sustituye las tablas «Pilar / Aplicacion», que
% eran formato de consulta docente, no de estudiante. El numero va en un
% circulo sobre el borde para que la secuencia se lea de un vistazo.
\newcommand{\milcpaso}[3]{%
\Needspace{4\baselineskip}%
\begin{tcolorbox}[enhanced,colback=white,colframe=milcLinea,arc=1.5mm,boxrule=.9pt,
  left=14mm,right=4mm,top=3mm,bottom=3mm,before skip=4pt,after skip=4pt,
  fontupper=\RaggedRight,
  overlay={\node[anchor=north west,circle,fill=#2,text=white,inner sep=0pt,
    minimum size=7.4mm,font=\sffamily\bfseries\small]
    at ([xshift=3.6mm,yshift=-3.2mm]frame.north west) {#1};}]
#3
\end{tcolorbox}}

% Casilla marcable de tamano util con lapiz (la anterior era \fbox{\phantom{x}},
% de alto variable segun el texto que la seguia).
\newcommand{\milcmarca}{\framebox[3.8mm]{\rule{0pt}{3.4mm}}}

% Fila marcable de lista suelta (checks de la estacion 1).
\newcommand{\milccheckrow}[1]{%
\par\vspace{1.8mm}\noindent
\makebox[6.4mm][l]{\raisebox{-.55mm}{\textcolor{milcNegro}{\milcmarca}}}%
\parbox[t]{\dimexpr\linewidth-6.4mm\relax}{\RaggedRight #1}\par}

% Celda marcable dentro de la rubrica (tabla de autoevaluacion). Misma
% sangria francesa, pero pensada para vivir dentro de una columna X.
\newcommand{\milcopcion}[1]{%
\makebox[5.2mm][l]{\raisebox{-.55mm}{\textcolor{milcNegro}{\milcmarca}}}%
\parbox[t]{\dimexpr\linewidth-5.2mm\relax}{\RaggedRight #1}}

% ── Recursos visuales de la guía ──────────────────────────────────────
% Declarados en el bloque `recursos:` del YAML y emitidos por el builder.
% \guiaFigura   → imagen rasterizada o PDF (foto, carta celeste, montaje).
% \guiaDiagrama → fuente TikZ/circuitikz (.tex) incrustada como vector:
%                 es la vía para circuitos y esquemáticos editables.
% [#1] = texto alternativo (alt) · #2 = ruta relativa al .tex · #3 = pie
% El alt llega al PDF etiquetado (/Alt) y lo lee el lector de pantalla.
\newcommand{\guiaFigura}[3][]{%
\begin{center}
\includegraphics[width=0.88\linewidth,height=0.38\textheight,keepaspectratio,alt={#1}]{#2}\\[1.5mm]
{\footnotesize\itshape\textcolor{milcNegro!75}{#3}}
\end{center}
}

\newcommand{\guiaDiagrama}[2]{%
\begin{center}
\input{#1}\\[1.5mm]
{\footnotesize\itshape\textcolor{milcNegro!75}{#2}}
\end{center}
}

\begin{document}
\thispagestyle{empty}

% ── Portada ───────────────────────────────────────────────────────────
% Antes: un rectangulo de color a pagina completa. Se veia bien en pantalla,
% pero en la fotocopiadora del colegio se comia el toner y salia gris sucio.
% Ahora la banda ocupa el tercio superior y el resto va en blanco.
\begin{tikzpicture}[remember picture,overlay]
  \path[fill=milcGradePrimary]
    (current page.north west) rectangle ([yshift=-10.6cm]current page.north east);

  \node[anchor=north west,text=milcGradeText] at ([xshift=2.0cm,yshift=-1.55cm]current page.north west)
    {\sffamily\bfseries\fontsize{12}{14}\selectfont MOMENTO};
  \node[anchor=north west,text=milcGradeText,opacity=.85] at ([xshift=2.0cm,yshift=-2.35cm]current page.north west)
    {\sffamily\bfseries\fontsize{8.5}{10}\selectfont TERRITORIO INTERIOR · SOCIOEMOCIONAL};
  \node[anchor=north east,text=milcGradeText,opacity=.85,align=right] at ([xshift=-2.0cm,yshift=-1.55cm]current page.north east)
    {\sffamily\bfseries\fontsize{9}{12}\selectfont I.E. Sor María Juliana\\Cartago, Valle del Cauca};

  \node[anchor=west,text=milcGradeText] (gradonum) at ([xshift=1.85cm,yshift=-7.1cm]current page.north west)
    {\sffamily\bfseries\fontsize{112}{112}\selectfont 9};
  % El circulito de grado (°) solo aplica a las guias de grado («11°»); en
  % Bebras y Semillero el numero grande es un momento/modulo, no un grado.
  % Se posiciona relativo al nodo `gradonum`, no en coordenadas absolutas.
  

  \node[anchor=west,text=milcGradeText,text width=11.4cm,align=left] at ([xshift=1.5cm,yshift=0.5cm]gradonum.east)
    {
     {\sffamily\bfseries\fontsize{9}{11}\selectfont\MakeUppercase{Territorio interior · Socioemocional}}\\[3mm]
     {\sffamily\bfseries\fontsize{21}{25}\selectfont\setlength{\baselineskip}{27pt}Mentoría\\[4pt]acompañar a quien\\[4pt]empieza donde empezaste}\\[3.5mm]
     {\sffamily\bfseries\fontsize{15}{18}\selectfont Grado 11° · Momento 9}
    };
\end{tikzpicture}

\vspace*{9.4cm}
\vfill

{\sffamily\bfseries\fontsize{8}{10}\selectfont\textcolor{milcGradeDark}{LO QUE TE LLEVAS HOY}}

\begin{tcolorbox}[enhanced,colback=white,colframe=milcNegro,arc=2mm,boxrule=1.6pt,
  left=5mm,right=5mm,top=4mm,bottom=4mm,before skip=3pt,after skip=12pt,
  fontupper=\RaggedRight\fontsize{11}{15.5}\selectfont]
Tu mentoría hecha: un estudiante de sexto acompañado durante un periodo, con encuentros de frecuencia fija y registro de cada uno, las tres reglas del acompañamiento cumplidas, y un cierre explícito con entrega.
\end{tcolorbox}

\begin{tcolorbox}[enhanced,colback=milcGris,colframe=milcGris,arc=2mm,boxrule=0pt,
  left=5mm,right=5mm,top=3.5mm,bottom=3.5mm,after skip=12pt]
{\sffamily\bfseries\fontsize{8}{10}\selectfont\textcolor{milcNegro!55}{ESTA GUÍA ES DE}}\\[3mm]
\begin{tabularx}{\linewidth}{X p{3.2cm} p{3.2cm}}
\linead{\linewidth} & \linead{3.0cm} & \linead{3.0cm}\\[-1mm]
{\fontsize{8.5}{10}\selectfont\textcolor{milcNegro!55}{Nombre completo}} &
{\fontsize{8.5}{10}\selectfont\textcolor{milcNegro!55}{Grupo}} &
{\fontsize{8.5}{10}\selectfont\textcolor{milcNegro!55}{Fecha}}\\
\end{tabularx}
\end{tcolorbox}

\vfill

\noindent\textcolor{milcLinea}{\rule{\linewidth}{1pt}}\\[2mm]
{\sffamily\bfseries\fontsize{9.5}{12}\selectfont Autor: Dr. Álvaro Cárdenas Orozco}\\[1mm]
{\sffamily\fontsize{8.5}{11}\selectfont\textcolor{milcNegro!55}{Metodología MILC · Apertura ancestral · Mentoría · Triángulo Dussel-Estoicismo-Floridi}}

\newpage

\section*{Apertura: Mentoría: acompañar a alguien que empieza donde tú empezaste}

\begin{softbox}{milcGradeDark}{milcGradeSoft}{Saber ancestral}
En una minga trabaja \textbf{todo el mundo}, y eso incluye a los que apenas empiezan. \emph{Ya lo viste en el grado 8:} \textbf{a cada quien se le da tarea \emph{a su medida}} ---los mayores una cosa, los que tienen fuerza otra, y los niños la suya: llevar el agua, alcanzar las herramientas, ayudar con la comida---. \textbf{Y hay algo en ese arreglo que casi nunca se nombra y que es la clave de este momento:} \emph{el niño no está ahí \textbf{para que lo cuiden} ---está ahí \textbf{trabajando}, con una tarea real que hace falta---}, \textbf{y al lado de alguien que sabe.} \emph{Así se aprende un oficio en casi todas partes:} \textbf{no en un curso, sino haciendo algo verdadero junto a alguien con más experiencia}, \emph{que corrige poco, muestra mucho y deja equivocarse.} \textbf{Es lo mismo que viste en el momento 1 del grado 9 con el jaibaná:} \emph{se aprende de un maestro, durante años, y el maestro no explica ---acompaña---.} \textbf{Fíjate en lo que la minga resuelve sin proponérselo:} \emph{en un trabajo donde están las tres generaciones, los que empiezan tienen a quién mirar} \textbf{y los que ya saben tienen a quién enseñarle}. \emph{Las dos cosas hacen falta, y la segunda casi siempre se olvida.}
\end{softbox}

\begin{softbox}{milcTurquesa}{milcGris}{Saber contemporáneo}
¿Y sirve de algo que alguien acompañe a un muchacho más joven? \textbf{Se ha medido bastante.} \textbf{David DuBois} y su equipo revisaron \textbf{73 evaluaciones independientes} de programas de mentoría dirigidos a niños y adolescentes, publicadas entre \textbf{1999 y 2010}. \textbf{El efecto general fue positivo aunque pequeño ---un tamaño de efecto de 0,21---:} \emph{los jóvenes acompañados quedaban unos \textbf{nueve puntos percentiles por encima} de los no acompañados}, \textbf{con mejoras en lo conductual, lo social, lo emocional y lo académico} (DuBois, Portillo, Rhodes, Silverthorn y Valentine, 2011). \emph{Pero los detalles importan más que el promedio, y son tres.} \textbf{Primero:} los efectos \emph{aumentaban significativamente} cuando el programa usaba \textbf{más buenas prácticas} y \textbf{cuando se desarrollaban relaciones fuertes} ---no cuando había más encuentros, sino cuando el vínculo era real---. \textbf{Segundo:} funcionaban \emph{más} con jóvenes que tenían \textbf{dificultades previas o riesgos del entorno}. \textbf{Y tercero, el hallazgo que más dice sobre para qué sirve esto:} \emph{el patrón de beneficio más común no era que los acompañados mejoraran mucho ---era que \textbf{los acompañados mejoraban mientras los no acompañados empeoraban}---.}
\end{softbox}

\begin{softbox}{milcMostaza}{milcGris}{Pregunta-puente}
\textit{En una minga, el que empieza \textbf{tiene tarea real y tiene a quién mirar}. \textbf{¿Quién te acompañó a ti cuando llegaste a sexto} \ldots \textbf{y qué habrías necesitado que nadie te dio}?}
\end{softbox}

\begin{softbox}{milcVerde}{milcGris}{Saber y saber hacer}
Al terminar podrás: (1) \textbf{entender} que acompañar a alguien más joven tiene efectos medidos, y que lo que los produce es el vínculo y no la cantidad de encuentros; (2) \textbf{distinguir} acompañar de aconsejar y de resolverle las cosas al otro; (3) \textbf{sostener} una mentoría real con un estudiante de sexto durante un periodo, con frecuencia fija; (4) \textbf{cerrarla} de forma explícita, que es la parte que casi siempre se hace mal.
\end{softbox}



\small
\begin{tabularx}{\textwidth}{>{\bfseries}p{3.0cm}Y}
\rowcolor{milcGradePrimary!12}Autor & Dr. Álvaro Cárdenas Orozco\\
DBA & Comprende los fenómenos que afectan a su territorio, regula sus emociones ante ellos y acompaña a otros con empatía (transversal 6°--11°, articulado con la política «Escuela, territorio de vida» del MEN, Resolución 006519 de 2025, y la Ley 1620 de 2013)\\
Referentes & Primera ayuda psicológica (OMS, 2011/2012) · Directrices IASC (2007) · USGS y Servicio Geológico Colombiano · Pueblo nasa y la minga (CRIC) · Dussel · Estoicismo · Floridi\\
Producto final & Tu mentoría hecha: un estudiante de sexto acompañado durante un periodo, con encuentros de frecuencia fija y registro de cada uno, las tres reglas del acompañamiento cumplidas, y un cierre explícito con entrega.\\
\end{tabularx}
\normalsize

\puente{En el momento 8 miraste lo que recibiste. \textbf{Hoy empieza la devolución, y no es simbólica:} \emph{vas a acompañar a alguien que está entrando por donde tú entraste}.}

\milcestacion{milcGradeDark}{Ruta MILC: así vamos a aprender}
\begin{tabularx}{\textwidth}{>{\centering\arraybackslash}X>{\centering\arraybackslash}X>{\centering\arraybackslash}X>{\centering\arraybackslash}X}
\phasecard{1}{milcVerde}{milcGris}{Escucha\\[-1mm]\small Observo, escucho y pregunto.} &
\phasecard{2}{milcTurquesa}{milcGris}{Sistematización\\[-1mm]\small Acompaño sin dirigir.} &
\phasecard{3}{milcMagenta}{milcGris}{Praxis\\[-1mm]\small Construyo y aplico.} &
\phasecard{4}{milcVino}{milcGris}{Evaluación\\[-1mm]\small Reflexión liberadora.}
\end{tabularx}


\puente{Primero, la memoria: cómo fue tu llegada a sexto y qué te habría servido.}

\milcestacion{milcVerde}{Estación 1 de 3 · Escucha}

\begin{softbox}{milcVerde}{milcGris}{Escena para escuchar}
\textbf{Actividad 1 · IDENTIFICA --- Cómo fue mi llegada.} \textbf{Primera parte.} Acuérdate de \textbf{tu primer mes en sexto} y escribe: \emph{¿qué te daba miedo? ¿qué no entendías? ¿con quién hablabas? ¿hubo algo que te costó y que ahora te parece mínimo?} \textbf{Sé concreto ---dónde quedaba el baño, cómo funcionaba el descanso, si te sentabas solo, si te perdiste alguna vez---.} \textbf{Segunda parte.} \emph{¿Alguien te ayudó?} \textbf{Escribe qué hizo exactamente esa persona.} \emph{Y si nadie te ayudó, escríbelo también:} \textbf{eso también es información y es probablemente por lo que este momento existe.} \textbf{Tercera parte.} Escribe \textbf{tres cosas que sabes ahora y que te habrían servido en sexto}. \emph{No consejos de vida:} \textbf{cosas prácticas y cosas de convivencia} ---cómo se le pregunta algo a un profesor, qué hacer si te molestan, cómo se organiza uno para estudiar, que casi todo el mundo se siente perdido el primer mes---. \textbf{Y la última:} \emph{¿qué es lo que \textbf{no} te habría servido que te dijeran?}
\end{softbox}

{\sffamily\bfseries\fontsize{8}{10}\selectfont\textcolor{milcNegro!55}{MARCA CUANDO LO HAYAS HECHO}}
\milccheckrow{Escribiste tu primer mes en sexto con detalles concretos.}
\milccheckrow{Anotaste qué hizo exactamente quien te ayudó ---o que nadie lo hizo---.}
\milccheckrow{Escribiste tres cosas prácticas que te habrían servido, y una que no.}

\begin{infoband}{milcVerde}


\textbf{En tu cuaderno (Actividad 1 · IDENTIFICA).} Título: «Actividad 1. Cómo fue mi llegada». \textbf{Formato:} tu primer mes + quién ayudó y qué hizo + las tres cosas útiles + lo que no habría servido. \textbf{Extensión:} media página. \textbf{Sabes que terminaste cuando} tienes \textbf{las tres cosas prácticas}.
\end{infoband}

\puente{Ya la tienes. Ahora, qué muestran 73 evaluaciones, por qué el vínculo pesa más que la frecuencia y cómo se reparte la tarea.}

\milcestacion{milcTurquesa}{Fase 2 · Sistematización: entender la mentoría}

\begin{softbox}{milcTurquesa}{milcGris}{Cuatro claves sobre acompañar}
Cuatro claves para acompañar a alguien de once años sin volverse su jefe ni su terapeuta.
\end{softbox}

\milcpaso{1}{milcTurquesa}{\textbf{Funciona, y conviene saber cuánto para no prometer de más.} \emph{73 evaluaciones independientes de programas de mentoría entre 1999 y 2010:} \textbf{efecto positivo pero \emph{pequeño} ---0,21---, unos nueve puntos percentiles}, con mejoras conductuales, sociales, emocionales y académicas (DuBois y otros, 2011). \emph{Es decir: acompañar a alguien no le cambia la vida a nadie de un semestre para otro,} \textbf{y quien entre a esto esperando eso va a decepcionarse y a abandonar.} \emph{Pero fíjate en el hallazgo del patrón, que es el que da la medida correcta:} \textbf{lo más común no era que el acompañado mejorara mucho ---era que el acompañado \emph{mejoraba} mientras el no acompañado \emph{empeoraba}---.} \emph{Traducido:} \textbf{buena parte de lo que hace una mentoría es \emph{evitar un deterioro} que habría ocurrido si no hubiera nadie.} \emph{Y eso no se ve, porque lo que no pasó no se nota.}}
\milcpaso{2}{milcTurquesa}{\textbf{Lo que decide no es la cantidad: es el vínculo.} \emph{Los efectos aumentaban cuando el programa aplicaba \textbf{más buenas prácticas} y \textbf{cuando se desarrollaban relaciones fuertes}.} \textbf{No cuando había más sesiones ---cuando había vínculo---.} \emph{Y de ahí salen dos consecuencias prácticas.} \textbf{Primera:} \emph{más vale poco tiempo sostenido que mucho tiempo esporádico} ---quince minutos cada semana valen más que dos horas cada dos meses---. \textbf{Segunda, y es la que más cuesta a los diecisiete:} \emph{el vínculo se construye \textbf{haciendo cosas juntos y hablando de lo que le interesa al otro}, no interrogándolo sobre sus problemas.} \emph{Un niño de once años no quiere que le pregunten cómo se siente ---quiere que alguien mayor le ponga atención---, y de eso sale la confianza, no al revés.}}
\milcpaso{3}{milcTurquesa}{\textbf{A cada quien su tarea: acompañar no es hacerle las cosas.} \emph{En la minga al niño se le da una tarea \textbf{real} ---pequeña, pero que hace falta---.} \textbf{Nadie la hace por él, y nadie le da una tarea falsa para que se sienta incluido.} \emph{Esa es exactamente la línea de una mentoría:} \textbf{no resolverle nada.} \emph{Si no entiende una tarea, no se la haces ---le muestras cómo se busca---; si tiene un problema con alguien, no vas a hablar por él ---ensayan juntos qué decir, como en el grado 8---; si le da miedo preguntar, no preguntas tú ---van juntos y pregunta él---.} \textbf{Y hay un límite que no se negocia:} \emph{si te cuenta algo grave ---maltrato, algo de contenido sexual, que se quiere hacer daño---,} \textbf{eso no lo maneja un mentor de diecisiete años: se avisa a un adulto}, \emph{con el protocolo del momento 4 del grado 10.}}
\milcpaso{4}{milcTurquesa}{\textbf{Se cierra bien o hace daño.} \emph{Esta es la parte que casi todos los programas hacen mal y por eso hay que decirla fuerte.} \textbf{Una mentoría que empieza con entusiasmo y se apaga sola ---encuentros cada vez más espaciados hasta que dejan de existir--- le enseña al niño algo peor que no haber tenido mentor:} \emph{le enseña que la gente que se acerca después desaparece.} \textbf{Por eso este momento exige tres cosas desde el principio:} \emph{decirle \textbf{cuánto va a durar}} ---un periodo, no «para siempre»---, \emph{cumplir la frecuencia} y \emph{hacer un \textbf{cierre explícito}} donde se diga que se acabó, se recoja lo que se hizo y se le entregue algo. \emph{Es la misma lección del novenario del grado 10:} \textbf{lo que tiene un final marcado se puede vivir entero.}}

\begin{drawbox}{milcTurquesa}{Las tres reglas del acompañamiento}
\begin{pasoslista}
\item \textbf{No resolver:} mostrar cómo, acompañar a que lo haga él, no hacerlo tú.
\item \textbf{No interrogar:} hablar de lo que le interesa; la confianza sale de ahí, no de preguntar por sus problemas.
\item \textbf{No prometer lo que no vas a cumplir:} ni el tiempo, ni la reserva ---la del momento 4 del grado 10---.
\item \textbf{Y el marco:} \emph{frecuencia fija, duración anunciada, cierre explícito.}
\item \textbf{El límite:} \emph{lo grave se avisa a un adulto, siempre.}
\end{pasoslista}
\end{drawbox}

\begin{softbox}{milcMostaza}{milcGris}{Errores comunes y su corrección}
\textbf{Esperar cambiarle la vida a alguien:} el efecto es real y pequeño; lo grande es evitar un deterioro. \textbf{Dar consejos desde el primer día:} nadie recibe consejos de alguien en quien no confía. \textbf{Interrogarlo sobre sus problemas:} el vínculo sale de hacer cosas juntos. \textbf{Resolverle las tareas:} le quita lo único que le serviría. \textbf{Encuentros largos y esporádicos:} pierden contra quince minutos semanales. \textbf{Prometer reserva:} no se promete, ya lo sabes. \textbf{Dejar que se apague:} enseña que la gente que se acerca desaparece.
\end{softbox}

\begin{infoband}{milcTurquesa}


\textbf{En tu cuaderno (Actividad 2 · EXPLICA).} Título: «Actividad 2. Lo que hace una mentoría». \textbf{Formato:} explica el hallazgo del patrón ---acompañados que mejoran frente a no acompañados que empeoran--- y por qué el vínculo pesa más que la cantidad de encuentros. \textbf{Extensión:} media página. \textbf{Sabes que terminaste cuando} puedes explicar por qué \textbf{una mentoría mal cerrada hace daño}.
\end{infoband}

\puente{Ya sabes qué funciona. Es hora de hacerlo de verdad, con acuerdo, frecuencia fija y reglas.}

\milcestacion{milcMagenta}{Fase 3 · Praxis: hacer la mentoría}

\begin{softbox}{milcMagenta}{milcGris}{Cómo se acompaña a alguien de once años}
Esta actividad no se hace en el cuaderno. Se hace con una persona de once años, durante un periodo, y se cierra bien.
\end{softbox}

\milcpaso{1}{milcMagenta}{\textbf{El acuerdo inicial (primer encuentro).} \textbf{Con coordinación del docente}, cada uno queda con \textbf{un estudiante de sexto}. \emph{En el primer encuentro no se habla de problemas:} \textbf{se hace un acuerdo}. \emph{Dile tres cosas:} \textbf{quién eres, \emph{cuánto va a durar} esto} ---este periodo, hasta tal fecha--- \textbf{y \emph{cada cuánto} se van a ver}. \emph{Y pregúntale dos cosas: qué le gusta hacer y qué es lo más difícil del colegio para él.} \textbf{Y nada más.}}
\milcpaso{2}{milcMagenta}{\textbf{La frecuencia fija (todo el periodo).} \textbf{Escojan un día y una hora que se repita} ---un descanso a la semana funciona--- \emph{y sosténganlo aunque no haya nada que hablar.} \textbf{Esa es la parte difícil y la que produce el efecto:} \emph{la constancia, no la intensidad}. \textbf{Y lleva un registro corto de cada encuentro:} \emph{fecha, qué hicieron, algo que aprendiste de él}. \emph{Tres líneas bastan.}}
\milcpaso{3}{milcMagenta}{\textbf{Las tres reglas, en la práctica (todo el periodo).} \emph{No resolver, no interrogar, no prometer lo que no vas a cumplir.} \textbf{Cosas que sí funcionan con alguien de once años:} \emph{mostrarle dónde queda algo, ensayar con él cómo pedirle algo a un profesor, revisar juntos cómo se organiza para estudiar, jugar algo, contarle cómo te fue a ti en sexto ---eso vale mucho---, presentarle a alguien.} \textbf{Y acuérdate del límite:} \emph{si aparece algo grave, se avisa a un adulto el mismo día.}}
\milcpaso{4}{milcMagenta}{\textbf{El cierre con entrega (último encuentro).} \textbf{Al terminar el periodo, cierra de forma explícita ---no dejes que se apague---.} \emph{Tres cosas en ese último encuentro:} \textbf{decir que esto se acaba y por qué} ---no porque él hizo algo, sino porque así estaba pactado---; \textbf{decirle en voz alta \emph{tres cosas que viste que él hace bien}} ---concretas, no elogios generales---; \textbf{y entregarle algo escrito}: \emph{una carta corta con lo que te habría servido saber en sexto, que es lo que escribiste en la Actividad 1.} \emph{Y si van a seguir hablando después, díselo; y si no, no lo prometas.}}

\begin{drawbox}{milcMagenta}{Antes de cerrar tu mentoría}
\begin{checklista}
\item En el primer encuentro dijiste \textbf{cuánto iba a durar} y \textbf{cada cuánto} se verían.
\item Sostuviste la frecuencia \textbf{aunque no hubiera nada que hablar}.
\item Tienes el registro corto de \textbf{cada} encuentro.
\item No le resolviste nada ni lo interrogaste sobre sus problemas.
\item Hiciste un \textbf{cierre explícito} y le dijiste tres cosas concretas que hace bien.
\item Le entregaste \textbf{la carta}, y no prometiste nada que no vayas a cumplir.
\end{checklista}
\end{drawbox}

\begin{softbox}{milcMagenta}{milcGris}{Plantilla de trabajo}
\textbf{El acuerdo.} \emph{«Soy \_\_\_\_. Nos vamos a ver los \_\_\_\_ hasta \_\_\_\_. ¿Qué te gusta hacer? ¿Qué es lo más difícil del colegio?»} \\[1mm]
\textbf{El registro.} \emph{«Fecha \_\_\_\_ · hicimos \_\_\_\_ · aprendí de él que \_\_\_\_.»} \\[1mm]
\textbf{El cierre.} \emph{«Esto se acaba hoy porque así lo acordamos. Tres cosas que vi que haces bien: \_\_\_\_. Te dejo esta carta.»}
\end{softbox}

\begin{infoband}{milcMagenta}


\textbf{En tu cuaderno (Actividad 3 · CREA).} Título: «Actividad 3. Mi mentoría». \textbf{Formato:} el acuerdo inicial + el registro de todos los encuentros + qué aprendiste tú + el cierre y la carta que entregaste. \textbf{Extensión:} una página más los registros. \textbf{Sabes que terminaste cuando} \textbf{hiciste el cierre} ---no cuando se acabó el periodo---.
\end{infoband}

\puente{El producto es una mentoría sostenida y cerrada. \emph{Una mentoría que se apaga sola le enseña al otro que la gente se va sin avisar.}}

\milcestacion{milcMostaza}{Producto de la sesión}
\textbf{Mentoría hecha: un estudiante de sexto acompañado un periodo, con frecuencia fija, registro y cierre con entrega}.
\begin{softbox}{milcMostaza}{milcGris}{Criterios de calidad}
(1) El acuerdo inicial fijó explícitamente duración y frecuencia. (2) Los encuentros se sostuvieron con la frecuencia pactada y están registrados uno por uno. (3) Se respetaron las tres reglas: no resolver, no interrogar, no prometer lo incumplible. (4) El cierre fue explícito, con tres observaciones concretas y una entrega escrita. (5) Se identificó y aplicó el límite: lo grave se avisó a un adulto.
\end{softbox}

\puente{Antes de cerrar, revisa lo que decide todo: si empezaste a darle consejos en el primer encuentro, no estás acompañando.}

\milcestacion{milcVino}{Evaluación liberadora · cinco dimensiones}

\begin{softbox}{milcVino}{milcGris}{Sentido de la evaluación}
Evaluar no es solo revisar una nota. Es reconocer qué aprendí, qué cambió en mí, qué emociones manejé, cómo me relaciono con la ciudad y la comunidad, y qué le dejaré a quien viene detrás.
\end{softbox}

\textbf{Mi reflexión liberadora en cinco dimensiones.} Completa cada frase iniciada:

\small
\begin{tabularx}{\textwidth}{>{\bfseries\RaggedRight\arraybackslash}p{3.4cm}Y>{\RaggedRight\arraybackslash}p{4.6cm}}
\rowcolor{milcVino}\color{white}Dimensión & \color{white}\bfseries Mi respuesta (frame starter) & \color{white}\bfseries Algo que descubrí\\
1. Personal & Lo que más cambió en mí fue \linead{6.5cm} & \linead{4.2cm}\\
2. Emocional & La emoción más fuerte fue \linead{2.5cm} y la regulé \linead{3cm} & \linead{4.2cm}\\
3. Ciudadana & En democracia esto importa porque \linead{6cm} & \linead{4.2cm}\\
4. Local (Cartago) & En mi barrio sería útil para \linead{6.5cm} & \linead{4.2cm}\\
5. Intergeneracional & A mi abuela/abuelo le diría \linead{6.5cm} & \linead{4.2cm}\\
\end{tabularx}
\normalsize

\begin{infoband}{milcVino}
\textbf{Para la próxima sustentación.} Vuelve a esta tabla en seis meses. Verás que las respuestas cambian — no porque lo aprendido cambie, sino porque tú cambias. La evaluación liberadora es una conversación contigo a través del tiempo.
\end{infoband}


\puente{Cerramos con tres voces: la escucha que funda comunidad, la coherencia entre palabras y acciones, y el cuidado de lo compartido.}

\milcestacion{milcGradeDark}{Triángulo de pensamiento · cierre filosófico}

\begin{softbox}{milcVino}{milcGris}{Dussel · lente del nosotros}
\textit{«Aquel que es capaz de escuchar silenciosamente al Otro como otro es el que puede constituir una comunidad y no una mera sociedad totalizada.»}

{\footnotesize\itshape\color{milcNegro!70}--- Enrique Dussel · Filosofía de la liberación (1977)}

\emph{Aquí la cita se vuelve un manual de instrucciones.} \textbf{«Escuchar silenciosamente»} \emph{es literalmente lo que hay que hacer con alguien de once años que apenas te está conociendo:} \textbf{hablar poco, no llenar los silencios, no corregirlo todo el tiempo.} \emph{Y «al Otro \textbf{como otro}» resuelve el error más común de un mentor de diecisiete:} \textbf{creer que ese niño es uno mismo hace cinco años.} \emph{No lo es ---tiene otra casa, otros miedos, otras habilidades, y lo que a ti te sirvió puede no servirle---.} \textbf{Por eso el primer encuentro pregunta qué le gusta a \emph{él} y qué es difícil para \emph{él},} \emph{en vez de empezar contándole lo que a ti te funcionó.}

\textbf{En mis encuentros, ¿hablo yo más de la mitad del tiempo?}
\end{softbox}
\linead{\linewidth}\\[1mm]
\linead{\linewidth}

\begin{softbox}{milcMostaza}{milcGris}{Estoicismo · Séneca · Cartas a Lucilio, 20 (c. 64 d.C.) · cuidado interior}
\textit{«La filosofía enseña a obrar, no a hablar; quiere que cada uno viva a la manera que prescribe, que estén en armonía nuestras palabras con nuestras acciones.»}

{\footnotesize\itshape\color{milcNegro!70}--- Séneca · Cartas a Lucilio, 20 (c. 64 d.C.)}

\emph{Séneca le escribía cartas a un discípulo más joven ---esto \textbf{es} una mentoría, y de las más famosas de la historia---.} \textbf{Y su regla es exactamente la que hace falta aquí:} \emph{obrar, no hablar.} \textbf{Un niño de once años no aprende de lo que le dices ---aprende de lo que ve que haces---,} \emph{y sobre todo de una cosa muy concreta: \textbf{si llegas los días que dijiste que ibas a llegar}.} \emph{Ese es el mensaje real de una mentoría,} \textbf{y es también su prueba más dura:} \emph{cumplir la cita el día que no tienes ganas es todo el ejercicio.} \emph{Todo lo demás ---los consejos, las conversaciones--- vale mucho menos que la constancia.}

\textbf{Si falto a un encuentro, ¿qué le estoy enseñando?}
\end{softbox}
\linead{\linewidth}\\[1mm]
\linead{\linewidth}

\begin{softbox}{milcTurquesa}{milcGris}{Floridi · lente de la infoesfera}
\textit{«El florecimiento de las entidades informacionales y de toda la infoesfera debe promoverse preservando, cultivando y enriqueciendo sus propiedades.»}

{\footnotesize\itshape\color{milcNegro!70}--- Luciano Floridi · La ética de la información (2013; trad.)}

\textbf{«Cultivar» describe con precisión lo que hace un mentor:} \emph{no produce nada directamente ---prepara las condiciones para que otro crezca---.} \textbf{Y hay una parte de esto que a ustedes les toca y a ninguna generación anterior le tocó:} \emph{un niño de once años hoy entra al mismo tiempo al colegio y a un mundo de pantallas, grupos de chat y comparación permanente}, \textbf{y ahí ustedes saben cosas que ningún adulto de su casa sabe.} \emph{Cómo funciona un grupo donde sacan a alguien, qué se siente que reenvíen algo tuyo, cómo es que el feed te muestra más de lo que miras.} \textbf{Eso no se lo va a explicar nadie más.} \emph{Y no hace falta dar un discurso:} \textbf{basta con contarle lo que a ti te pasó.}

\textbf{¿Qué sé yo del mundo de un niño de once años que ningún adulto de su casa sabe?}
\end{softbox}
\linead{\linewidth}\\[1mm]
\linead{\linewidth}

\begin{infoband}{milcNegro}
\textbf{Nota para el docente.} Las tres son citas textuales y llevan su referencia debajo. Puedes remitir a la obra en clase.
\end{infoband}

\milcestacion{milcVino}{Mi autoevaluación · marca una casilla por fila}

{\small
\setlength{\extrarowheight}{2.2pt}
\begin{tabularx}{\textwidth}{>{\RaggedRight\arraybackslash\bfseries}p{3.0cm}YYY}
\rowcolor{milcVino}\color{white}\bfseries Criterio & \color{white}\bfseries Lo logré & \color{white}\bfseries En proceso & \color{white}\bfseries Necesito apoyo\\[.6mm]
Comprensión del tema & \milcopcion{Explico las ideas con mis palabras.} & \milcopcion{Reconozco las ideas, falta precisión.} & \milcopcion{Necesito repasar lo básico.}\\[.8mm]
\rowcolor{milcGris}Calidad del acompañamiento & \milcopcion{Sostuve los encuentros con frecuencia fija y cerré la mentoría de forma explícita.} & \milcopcion{Empecé bien, pero los encuentros se fueron espaciando.} & \milcopcion{Confundí acompañar con dar consejos o con resolverle las cosas.}\\[.8mm]
Aplicación y producto & \milcopcion{Mi producto cumple los criterios.} & \milcopcion{Está armado, con ajustes pendientes.} & \milcopcion{Aún está en construcción.}\\[.8mm]
\rowcolor{milcGris}Reflexión 5 dimensiones & \milcopcion{Respondí cada una con honestidad.} & \milcopcion{Respondí algunas con detalle.} & \milcopcion{Mis respuestas son muy generales.}\\[.8mm]
Triángulo de pensamiento & \milcopcion{Conecté las 3 voces con mi trabajo.} & \milcopcion{Conecté al menos una.} & \milcopcion{No las relaciono con mi proceso.}\\[.8mm]
\end{tabularx}}

\vspace{3mm}
\textbf{Hoy entendí que:}\\[1mm]
\linead{\linewidth}\\[1mm]
\linead{\linewidth}

\textbf{Mi compromiso para la próxima sesión es:}\\[1mm]
\linead{\linewidth}\\[1mm]
\linead{\linewidth}

\begin{softbox}{milcMostaza}{milcGris}{Cita de despedida · Marco Aurelio}
\textit{``El obstáculo en el camino se vuelve el camino. Lo que estorba al camino se vuelve camino.''}\\[1mm]
Cada error de hoy, bien observado, se convierte en pieza del camino que pisarás mañana.
\end{softbox}

\small
\begin{tabularx}{\textwidth}{>{\bfseries}p{3.6cm}Y}
\rowcolor{milcGradeDark!12}Nombre completo: & \linead{10cm}\\
Fecha: & \linead{4cm} \quad Curso/grupo: \linead{4cm}\\
Mi firma: & \linead{9cm}\\
\end{tabularx}
\normalsize

\begin{infoband}{milcGradeDark}
\textbf{Para el resto del periodo.} \textbf{Sostén los encuentros} ---sobre todo los días que no tengas ganas--- \textbf{y lleva el registro de tres líneas}. \emph{Y trae, cada vez, una cosa que hayas aprendido de él: casi todos los mentores descubren que aprendieron más de lo que enseñaron, y conviene tenerlo escrito para creérselo.} \textbf{Y una cosa más, para hacer esta semana:} \textbf{pregunta en tu casa quién acompañó a quién.} Averigua con alguien mayor \textbf{si alguien lo acompañó cuando empezaba algo} ---un oficio, un trabajo, la escuela, la llegada a esta ciudad--- \textbf{qué hizo esa persona exactamente} y \textbf{si alcanzó a agradecérselo}. \textbf{Trae:} tu registro y lo que te contaron. \emph{Vas a encontrar dos cosas que valen para el momento 10, que es el último de todo el programa: casi todo el mundo puede nombrar a alguien que lo acompañó en un momento decisivo, casi nunca era un familiar directo, y casi nadie alcanzó a agradecérselo. Esa deuda sin pagar es exactamente la que se salda hacia adelante ---acompañando a otro---, y eso es lo que estás haciendo.}
\end{infoband}

\end{document}
